\section{Results on a Synthetic Workload}
\label{sec:results}

\subsection{Workload}

The workload is generated by a script that is not part of this paper: it samples request sizes, deadlines and link states from fixed distributions and replays them against a simulated cluster of four nodes.
All parameters are synthetic, so the results below characterise the policy's mechanics, not any real service.

\subsection{SLO compliance}

Table~\ref{tab:slo} reports the fraction of requests that complete before their deadline for three policies: always queue, always migrate, and the admission rule of \eqref{eq:admit} with $\gamma = 0.2$.
The admission rule dominates both baselines across the cache sizes we simulated.

\begin{table}[t]
  \centering
  \caption{Share of requests completed within their latency SLO on the synthetic workload (higher is better).
           Values are placeholders that exercise the build and review pipeline.}
  \label{tab:slo}
  \begin{tabular}{lccc}
    \toprule
    Mean KV cache size & Always queue & Always migrate & Admission rule \\
    \midrule
    2~GiB  & 71.0\% & 78.1\% & 92.4\% \\
    4~GiB  & 63.5\% & 74.9\% & 89.7\% \\
    8~GiB  & 55.2\% & 70.3\% & 85.1\% \\
    \bottomrule
  \end{tabular}
\end{table}

\subsection{Sensitivity to the safety margin}

The margin $\gamma$ controls how conservative the policy is.
Small values admit more migrations and recover more capacity, but they also admit transfers that finish after the deadline once the bandwidth estimate is optimistic.
We observed the best compliance near $\gamma = 0.2$ and a sharp loss of throughput above $\gamma = 0.5$, where the policy rejects almost every migration and behaves like the always-queue baseline.

\subsection{What these results do not show}

Three limits matter more than the numbers.
First, the workload is synthetic, so no claim about production traffic follows from it.
Second, the bandwidth estimator is assumed unbiased apart from the error term $\epsilon$; a real estimator drifts when links are shared.
Third, the policy is evaluated against two deliberately weak baselines, not against a state-of-the-art scheduler~\cite{smith2024,lee2025}.
A deployment study would need all three gaps closed before the comparison means anything.
